\documentclass[11pt]{article}

\usepackage[utf8]{inputenc}
\usepackage{amsmath}
\usepackage{amssymb}
\usepackage{algorithm}
\usepackage{algpseudocode}
\usepackage[margin=2.5cm]{geometry}

\title{The $p$-Median Problem:\\ Model and Greedy Heuristic}
\author{}
\date{}

\begin{document}
\maketitle

\section{The $p$-median model}

We are given a set of $n$ demand points (hexes) and the same set is used as the
set of candidate facility sites (medians). Let

\begin{itemize}
  \item $V = \{1, \dots, n\}$ be the set of hexes (both demands and candidate sites);
  \item $d_{ij} \ge 0$ be the distance between hex $i$ and hex $j$;
  \item $p$ be the number of facilities (open sites) to select.
\end{itemize}

In this implementation the distance is the Euclidean distance between hex grid
locations,
\[
  d_{ij} = \sqrt{(x_i - x_j)^2 + (y_i - y_j)^2},
\]
where $(x_i, y_i)$ is the location of hex $i$.

\subsection*{Decision variables}

\[
  x_{ij} \quad\text{the fraction of demand } i \text{ served by site } j,
  \qquad i, j \in V.
\]

There is no separate facility-open variable. The role of $y_j$ is played by the
\emph{diagonal} entry $x_{jj}$: a hex $j$ serves itself exactly when it is open,
so $x_{jj}$ is the open indicator and is the only variable that must be integral,
\[
  x_{jj} \in \{0, 1\} \qquad j \in V,
  \qquad\qquad
  0 \le x_{ij} \le 1 \quad (i \ne j).
\]
The off-diagonal assignments $x_{ij}$ ($i \ne j$) stay continuous in $[0,1]$.
Once the $p$ open sites are fixed by the binary diagonal, the remaining
assignment problem is totally unimodular, so the continuous $x_{ij}$ take values
in $\{0,1\}$ at the optimum anyway.

\subsection*{Formulation}

\begin{align}
  \min_{x} \quad
    & \sum_{i \in V} \sum_{j \in V} d_{ij}\, x_{ij}
      \label{eq:obj}\\
  \text{s.t.} \quad
    & \sum_{j \in V} x_{ij} = 1
      & \forall\, i \in V, \label{eq:assign}\\
    & x_{ij} \le x_{jj}
      & \forall\, i, j \in V, \label{eq:link}\\
    & \sum_{j \in V} x_{jj} = p,
      & \label{eq:count}\\
    & x_{jj} \in \{0, 1\}
      & \forall\, j \in V, \label{eq:bin}\\
    & 0 \le x_{ij} \le 1
      & \forall\, i \ne j. \label{eq:dom}
\end{align}

The objective~\eqref{eq:obj} minimises the total assignment distance.
Constraint~\eqref{eq:assign} forces every hex to be fully assigned,
\eqref{eq:link} allows a hex to be served by site $j$ only if $j$ is open
(i.e.\ $x_{jj} = 1$), and \eqref{eq:count} opens exactly $p$ facilities.

\subsection*{Weighted demand}

Each hex carries a demand weight $w_h \ge 0$ (in the code, \texttt{Hex::weight},
default $1$). The weighted $p$-median replaces the objective~\eqref{eq:obj} with
\[
  \min_{x} \quad \sum_{i \in V} \sum_{j \in V} w_i\, d_{ij}\, x_{ij},
\]
so that hexes with larger demand pull facilities closer. The weights also drive
the workload-balancing repair pass below, where the \emph{load} of a district is
the sum of the weights of the hexes assigned to it.

\section{Greedy heuristic}

The greedy heuristic builds the solution by repeatedly opening the candidate
site that reduces the total assignment distance the most. Let $c_h$ be the
distance from hex $h$ to its nearest currently open site; opening site $c$ would
change the total cost to $\sum_{h \in V} \min(c_h, d_{hc})$.

\begin{algorithm}[h]
\caption{Greedy $p$-median (site selection)}
\label{alg:greedy}
\begin{algorithmic}[1]
  \Require hexes $V$, distances $d$, facility count $p$
  \Ensure set of open sites $S$
  \State $c_h \gets \infty$ \textbf{ for all } $h \in V$
      \Comment{cost of each hex to its nearest open site}
  \State $S \gets \emptyset$
  \For{$k \gets 1$ \textbf{ to } $p$}
    \State $c^\star \gets \displaystyle\arg\min_{c \in V \setminus S}
            \sum_{h \in V} \min\!\big(c_h,\, d_{hc}\big)$
        \Comment{site giving the lowest total cost}
    \State $S \gets S \cup \{c^\star\}$
    \For{$h \in V$}
      \State $c_h \gets \min(c_h,\, d_{h c^\star})$
          \Comment{update nearest-site costs}
    \EndFor
  \EndFor
  \State \Return $S$
\end{algorithmic}
\end{algorithm}

Once the $p$ sites are chosen, each hex is assigned to its nearest open site:
\[
  \text{assign}(h) = \arg\min_{s \in S} d_{hs}, \qquad h \in V,
\]
and the objective is $\sum_{h \in V} d_{h,\,\text{assign}(h)}$.

\begin{algorithm}[h]
\caption{Nearest-site assignment}
\label{alg:assign}
\begin{algorithmic}[1]
  \Require hexes $V$, distances $d$, open sites $S$
  \Ensure assignment $a$
  \For{$h \in V$}
    \State $a_h \gets \displaystyle\arg\min_{s \in S} d_{hs}$
  \EndFor
  \State \Return $a$
\end{algorithmic}
\end{algorithm}

The greedy solution may then be passed through optional repair passes
(contiguity and workload balancing) before reporting the final objective.

\section{Repair passes}

The repair passes operate in \emph{district-index space}: $a_h \in \{1,\dots,p\}$
is the index of the open site currently serving hex $h$. They rely on the hex
adjacency graph $N(h)$, the set of grid neighbours of $h$ (the eight
king-move offsets on the hex grid).

\subsection*{Contiguity repair}

A district should be a single connected region. For each district $s$, a
breadth-first search from its open site $S_s$ marks every hex of $s$ reachable
through neighbours that also belong to $s$. Any hex of $s$ that is \emph{not}
reachable is ``stranded'' and is reassigned to a neighbouring district.

\begin{algorithm}[h]
\caption{Contiguity repair}
\label{alg:contig}
\begin{algorithmic}[1]
  \Require assignment $a$, adjacency $N$, open sites $S$
  \Ensure repaired assignment $a$
  \For{$s \gets 1$ \textbf{ to } $p$}
    \State $R \gets \textsc{Bfs}(S_s,\, N,\, a,\, s)$
        \Comment{hexes of $s$ reachable from its site}
    \For{$h \in V$ \textbf{ with } $a_h = s$ \textbf{ and } $h \notin R$}
      \If{$\exists\, j \in N(h)$ \textbf{ with } $a_j \ne s$}
        \State $a_h \gets a_j$
            \Comment{move stranded hex to a neighbour's district}
      \EndIf
    \EndFor
  \EndFor
  \State \Return $a$
\end{algorithmic}
\end{algorithm}

Here $\textsc{Bfs}(\text{start}, N, a, s)$ returns the set of hexes reachable
from \emph{start} using only edges of $N$ whose endpoints satisfy $a_j = s$.

\subsection*{Workload-balancing repair}

This pass evens out district loads until the most loaded district is within a
factor \texttt{limit} of the mean load. Let
$\mathrm{load}(s) = \sum_{h : a_h = s} w_h$ and
$\overline{\mathrm{load}} = \frac{1}{p}\sum_s \mathrm{load}(s)$. While the
imbalance $\max_s \mathrm{load}(s) / \overline{\mathrm{load}}$ exceeds the limit,
a single boundary hex is moved from the most loaded district to the least loaded
one, provided the move keeps the donor district connected.

\begin{algorithm}[h]
\caption{Workload-balancing repair}
\label{alg:workload}
\begin{algorithmic}[1]
  \Require assignment $a$, weights $w$, adjacency $N$, open sites $S$, ratio \textit{limit}
  \Ensure repaired assignment $a$
  \For{$\textit{iter} \gets 1$ \textbf{ to } $\textit{maxIter}$}
    \State compute $\mathrm{load}(s)$ for all $s$, and mean $\overline{\mathrm{load}}$
    \If{$\overline{\mathrm{load}} = 0$ \textbf{ or } $\max_s \mathrm{load}(s) / \overline{\mathrm{load}} \le \textit{limit}$}
      \State \textbf{break}
    \EndIf
    \State $o \gets \arg\max_s \mathrm{load}(s)$;\quad $u \gets \arg\min_s \mathrm{load}(s)$
    \If{$o = u$} \State \textbf{break} \EndIf
    \State \textit{swapped} $\gets$ \textbf{false}
    \For{$h \in V$ \textbf{ with } $a_h = o$ \textbf{ and } $h \ne S_o$}
      \If{$\exists\, j \in N(h)$ \textbf{ with } $a_j = u$ \textbf{ and } \textsc{ConnectedWithout}$(a, N, S_o, h, o)$}
        \State $a_h \gets u$;\quad \textit{swapped} $\gets$ \textbf{true};\quad \textbf{break}
      \EndIf
    \EndFor
    \If{\textbf{not} \textit{swapped}} \State \textbf{break} \EndIf
  \EndFor
  \State \Return $a$
\end{algorithmic}
\end{algorithm}

\textsc{ConnectedWithout}$(a, N, S_o, h, o)$ tests, via a BFS from the site
$S_o$ that skips hex $h$, whether district $o$ would remain connected after $h$
is removed (i.e.\ the BFS reaches all but one of the district's hexes). The loop
also stops if no admissible move exists or after \textit{maxIter} iterations,
guaranteeing termination.

\end{document}
